\begin{helper}
Let $V$ be finite and let $\nu$ be a measure on
$\mathcal P(V)\times\mathbb R^V$, with the discrete measurable structure
on the first factor and the Borel structure on the second.
Assume each activation fibre $H_A=\{\eta:\eta_{\rm act}=A\}$ has finite
mass, and for every $A\subseteq V$ and $t\in[0,1]^V$ assume
\[
 \nu\{\eta\in H_A:0\le\pi(v)\le t_v\ (v\in V)\}
 =\nu(H_A)\operatorname{ofReal}\Bigl(\prod_{v\in V}t_v\Bigr).
\]
Then for $\nu$-almost every sample, $0<\pi(v)\le1$ for every $v\in V$.
\end{helper}
\begin{proof}
Put $U=\{\eta:0\le\pi(v)\le1\text{ for all }v\in V\}$.
All activation fibres and coordinate interval events are measurable.
The rectangle formula at $t\equiv1$ gives
$\nu(H_A\cap U)=\nu(H_A)$. Since this mass is finite, subtraction for
the nested measurable sets gives $\nu(H_A\setminus U)=0$.
The finitely many activation fibres cover the whole space; hence
$\nu(U^c)=0$ by finite subadditivity.
For any fixed $v_0\in V$, use the threshold $t_{v_0}=0$ and $t_v=1$
for $v\ne v_0$. Its product is zero. The associated rectangle on $H_A$
is exactly $H_A\cap U\cap\{\pi(v_0)=0\}$, and therefore is null.
Taking the finite union over $A$ shows $U\cap\{\pi(v_0)=0\}$ is null.
Taking the finite union over $v_0$, and adding $U^c$, gives a null set
outside which every coordinate is in $[0,1]$ and is not zero.
Thus every coordinate lies in $(0,1]$, as claimed. If $V$ is empty,
the conclusion is vacuous and the same argument has no boundary terms.
\end{proof}
